\documentclass[11pt]{article}
\usepackage[T1]{fontenc}
\usepackage[utf8]{inputenc}
\usepackage{lmodern}
\usepackage[letterpaper,margin=0.78in]{geometry}
\usepackage{microtype}
\usepackage{booktabs,longtable,tabularx,array,multirow}
\usepackage{amsmath,amssymb}
\usepackage{xcolor}
\usepackage{enumitem}
\usepackage{titlesec}
\usepackage{fancyhdr}
\usepackage{lastpage}
\usepackage{hyperref}
\usepackage[most]{tcolorbox}
\usepackage{parskip}
\usepackage{needspace}
\usepackage{ragged2e}
\usepackage{setspace}

\definecolor{Navy}{HTML}{17365D}
\definecolor{Blue}{HTML}{2F75B5}
\definecolor{Teal}{HTML}{1F6D6A}
\definecolor{Gold}{HTML}{B38B2D}
\definecolor{Gray}{HTML}{5E6B75}
\definecolor{LightBlue}{HTML}{EEF4FA}
\definecolor{LightGold}{HTML}{FBF6E8}
\definecolor{LightGray}{HTML}{F4F5F7}
\definecolor{LightTeal}{HTML}{EDF7F6}

\hypersetup{colorlinks=true,linkcolor=Navy,urlcolor=Blue,citecolor=Teal,pdfauthor={OpenAI},pdfsubject={Value investing study and implementation materials}}
\setlist[itemize]{leftmargin=1.5em,itemsep=2pt,topsep=3pt}
\setlist[enumerate]{leftmargin=1.7em,itemsep=2pt,topsep=3pt}
\renewcommand{\arraystretch}{1.16}
\setlength{\tabcolsep}{5pt}
\titleformat{\section}{\Large\bfseries\color{Navy}}{\thesection}{0.6em}{}[\vspace{-0.25em}\color{Blue}\rule{\linewidth}{0.5pt}]
\titleformat{\subsection}{\large\bfseries\color{Teal}}{\thesubsection}{0.5em}{}
\titleformat{\subsubsection}{\normalsize\bfseries\color{Gray}}{\thesubsubsection}{0.5em}{}
\pagestyle{fancy}
\fancyhf{}
\fancyhead[L]{\small\color{Gray}\DocShortTitle}
\fancyhead[R]{\small\color{Gray}Value Investing Reference Set}
\fancyfoot[C]{\small\color{Gray}Page \thepage\ of \pageref*{LastPage}}
\renewcommand{\headrulewidth}{0.4pt}
\renewcommand{\footrulewidth}{0pt}

\newtcolorbox{principlebox}[1][]{enhanced,breakable,colback=LightBlue,colframe=Blue,boxrule=0.6pt,arc=2mm,left=2mm,right=2mm,top=1.5mm,bottom=1.5mm,title={#1},fonttitle=\bfseries}
\newtcolorbox{cautionbox}[1][]{enhanced,breakable,colback=LightGold,colframe=Gold,boxrule=0.6pt,arc=2mm,left=2mm,right=2mm,top=1.5mm,bottom=1.5mm,title={#1},fonttitle=\bfseries}
\newtcolorbox{modernbox}[1][]{enhanced,breakable,colback=LightTeal,colframe=Teal,boxrule=0.6pt,arc=2mm,left=2mm,right=2mm,top=1.5mm,bottom=1.5mm,title={#1},fonttitle=\bfseries}
\newtcolorbox{worksheetbox}[1][]{enhanced,breakable,colback=white,colframe=Gray,boxrule=0.45pt,arc=1mm,left=2mm,right=2mm,top=1.5mm,bottom=1.5mm,title={#1},fonttitle=\bfseries}
\newcommand{\checkbox}{\(\square\)}
\newcommand{\blankline}{\rule{0.96\linewidth}{0.35pt}}
\newcommand{\sourceRef}[1]{\textit{Source anchor: #1}}
\newcolumntype{Y}{>{\RaggedRight\arraybackslash}X}
\newcolumntype{C}[1]{>{\centering\arraybackslash}p{#1}}

\newcommand{\makedoccover}[4]{%
\begin{titlepage}
\thispagestyle{empty}
\vspace*{1.0in}
{\color{Navy}\rule{\linewidth}{2.2pt}}\\[0.55in]
{\Huge\bfseries\color{Navy} #1\par}
\vspace{0.18in}
{\Large\color{Teal} #2\par}
\vspace{0.5in}
\begin{tcolorbox}[colback=LightGray,colframe=Gray,boxrule=0.5pt,arc=2mm]
\textbf{Purpose.} #3
\end{tcolorbox}
\vfill
{\large Prepared from the revised edition of \textit{The Intelligent Investor}, with Jason Zweig's commentary.\par}
\vspace{0.12in}
{\small #4\par}
\vspace{0.25in}
{\small\color{Gray}\parbox{0.90\textwidth}{\centering Educational material only. Historical thresholds and examples are preserved for study; they are not individualized investment, tax, or legal advice.}}\par
\vspace{0.35in}
{\color{Navy}\rule{\linewidth}{2.2pt}}
\end{titlepage}
}

\newcommand{\DocShortTitle}{Modern Framework}
\begin{document}
\makedoccover{Modern Self-Directed Brokerage and\\Long-Term Portfolio Framework}{A source-anchored implementation for contemporary self-directed investing}{Maps the durable principles to account architecture, automation, execution discipline, diversification, tax-aware process, and ongoing governance.}{Standalone LaTeX study document}
\tableofcontents\newpage

\begin{cautionbox}[Source and modernization boundary]
The underlying source is the fourth revised edition, with Benjamin Graham's text updated in 1971--1972 and later commentary by Jason Zweig. Historical yields, tax rules, company examples, market levels, and dollar thresholds are therefore period-specific. This document preserves the source's framework and labels later implementation ideas separately rather than silently treating historical figures as current facts.
\end{cautionbox}

\section{Purpose and Design Principles}
This document translates the source's durable principles into a modern self-directed workflow without pretending that its 1971--1972 market data or 2003 commentary are current rules. The implementation uses six controls: low structural cost, broad diversification, explicit asset allocation, automated contributions, valuation discipline for any direct-stock sleeve, and a hard separation between investment and speculation.

\section{Account Architecture}
\begin{modernbox}[Conceptual account stack]
\textbf{Liquidity layer:} emergency cash and near-term spending reserves outside the risk portfolio.\\
\textbf{Long-term core:} diversified equity and high-quality fixed-income exposure implemented with low-cost funds or carefully selected securities.\\
\textbf{Enterprising sleeve:} optional direct securities that pass a written research and margin-of-safety process.\\
\textbf{Speculative sleeve:} optional, capped, separately tracked, and never necessary for the plan to succeed.
\end{modernbox}

A taxable brokerage account is useful for long-horizon assets that do not fit or are not intended for tax-advantaged accounts. Tax treatment is jurisdiction- and year-specific; the source's tax appendix is historical and should not be used as a current tax rulebook.

\section{Portfolio Policy Framework}
\subsection{Source-anchored allocation}
The book's defensive reference is a 50/50 stock-bond mix, bounded by 25\% and 75\% in either major asset class. A modern investor can treat that as a governance framework rather than as a universal target:
\begin{itemize}
\item choose a strategic equity target inside a range you can maintain through severe drawdowns;
\item hold the remainder in high-quality fixed income consistent with liquidity and duration needs;
\item define the rebalance rule in advance;
\item alter strategic allocation only when horizon, liquidity, or capacity for loss changes materially.
\end{itemize}

\subsection{Illustrative Graham-anchored reference portfolio}
This example is educational, not individualized advice:
\begin{center}
\begin{tabularx}{0.84\linewidth}{@{}Y C{0.18\linewidth}@{}}
\toprule
\textbf{Sleeve} & \textbf{Illustrative weight}\\
\midrule
Broad diversified equities & 50\%\\
High-quality nominal fixed income & 40\%\\
Inflation-protected fixed income & 10\%\\
\bottomrule
\end{tabularx}
\end{center}
The equity portion may itself be diversified across domestic and international markets. The later commentary suggests that holding part of stock exposure outside the investor's home market can reduce single-country concentration.

\section{Core-Plus-Satellite Implementation}
\subsection{Core}
The defensive core should be broad, inexpensive, and difficult to misuse. A total-market or similarly diversified index approach is consistent with the commentary's recommendation for defensive investors because it minimizes security-selection burden and turnover.

\subsection{Enterprising satellite}
If direct security selection is used, fund it by carving capital from the equity allocation rather than by adding unbounded risk on top of the policy. Every holding should have:
\begin{enumerate}
\item a business-quality assessment;
\item normalized financial statements;
\item a conservative valuation range;
\item an explicit margin of safety;
\item a maximum position size;
\item written disconfirming evidence and review triggers.
\end{enumerate}

\section{Funding and Automation}
\subsection{Contribution sequence}
\begin{enumerate}
\item Satisfy near-term liquidity needs first.
\item Direct recurring contributions to the most underweight long-term sleeve.
\item Use automatic investment where available to reduce timing decisions.
\item Rebalance with cash flows before selling appreciated assets in taxable accounts.
\end{enumerate}

\subsection{Dollar-cost averaging}
Periodic investing is valuable primarily as a behavioral and cash-flow discipline. It does not guarantee a profit and should not be confused with a claim that gradually investing a lump sum always dominates immediate investment.

\section{Order-Execution Discipline}
The source is not an order-type manual, but its price/value discipline translates naturally into execution controls.
\begin{center}
\small
\begin{tabularx}{\linewidth}{@{}p{0.22\linewidth}Y Y@{}}
\toprule
\textbf{Order behavior} & \textbf{Use case} & \textbf{Value-investing control}\\
\midrule
Market order & Highly liquid security where immediate execution matters more than a small price difference. & Confirm the spread and avoid using it in thin or volatile securities without understanding slippage.\\
Limit order & Direct-stock purchase where the maximum acceptable price is known. & Set the limit from valuation discipline, not from arbitrary round numbers.\\
Stop / stop-limit & Risk-management tool for trading contexts. & Do not substitute a price trigger for business-value analysis in a long-term investment thesis.\\
Recurring purchase & Core diversified funds. & Supports automatic policy implementation and reduces market-timing behavior.\\
\bottomrule
\end{tabularx}
\end{center}

\section{Margin and Leverage}
The book's warnings about speculation and the commentary's explanation of margin support a strong default: a long-term self-directed investment plan should not require borrowed money to meet its return objective. Leverage can convert ordinary quotation volatility into forced selling, defeating the ability to wait for value to be recognized.

\section{Tax-Aware Process Without Historical Tax Rates}
\begin{itemize}
\item Track cost basis and holding periods accurately.
\item Consider tax consequences before rebalancing through sales in taxable accounts.
\item Prefer contribution-directed rebalancing when practical.
\item Evaluate fund turnover and distribution characteristics where relevant.
\item Do not use the source's 1972 or 2003 tax rules as current law; verify current rules separately.
\end{itemize}

\section{Modern Risk Controls}
\subsection{Concentration}
Diversification is the default protection against being wrong. Direct-stock positions should be sized so that a permanent impairment does not derail the entire plan.

\subsection{Sequence and liquidity risk}
Keep short-horizon obligations out of volatile assets. Market prices can remain depressed longer than expected; the ability to avoid forced selling is itself a margin of safety.

\subsection{Behavioral risk}
Set rules for media consumption, portfolio-check frequency, and trade authorization. If a decision is triggered only by a quote, headline, or social proof, require a cooling-off interval and a written valuation update before execution.

\section{A Modernized Security-Research Packet}
For any direct equity holding, maintain one page for each of these modules:
\begin{enumerate}
\item business model and competitive economics;
\item five-to-ten-year financial summary;
\item debt, liquidity, and maturity profile;
\item normalized EPS and free-cash-flow bridge;
\item management and capital-allocation record;
\item valuation triangulation;
\item margin-of-safety assessment;
\item thesis risks and disconfirming evidence;
\item position sizing and exit/review policy.
\end{enumerate}
This packet operationalizes Graham's insistence on analysis before price action and Chapter 12's skepticism toward headline earnings.

\section{Rebalancing Policy}
A simple policy might state:
\begin{principlebox}[Illustrative rule]
Review allocation quarterly. Rebalance with contributions whenever possible. Execute sales only when a sleeve breaches its pre-defined tolerance band, when risk capacity changes, or when a security thesis materially deteriorates. Do not rebalance because of a forecast about next month's market direction.
\end{principlebox}

\section{Performance Measurement}
Measure success against the policy's objective rather than against the best-performing asset of the year. Track:
\begin{itemize}
\item time-weighted portfolio return versus an appropriate blended benchmark;
\item money-weighted return to understand the impact of contributions and behavior;
\item all-in fees and tax drag;
\item turnover;
\item maximum drawdown and recovery time;
\item whether written decision rules were followed.
\end{itemize}
A process can be successful even when it temporarily underperforms; an undisciplined speculative process can be unsuccessful even after a profitable year.

\section{Operating Calendar}
\begin{center}
\begin{tabularx}{\linewidth}{@{}p{0.20\linewidth}Y@{}}
\toprule
\textbf{Frequency} & \textbf{Action}\\
\midrule
Pay period / monthly & Contribute automatically; direct cash toward underweights.\\
Quarterly & Check allocation, liquidity, and major thesis-break events.\\
After earnings / filings & For direct holdings, update normalized earnings and balance-sheet risk only when new fundamental information warrants it.\\
Annually & Review policy, fees, diversification, tax records, adviser/broker settings, beneficiaries, and speculative-cap compliance.\\
After major life change & Rebuild liquidity and allocation assumptions before changing security selection.\\
\bottomrule
\end{tabularx}
\end{center}

\section{Sample Investment Policy Language}
\begin{modernbox}[Policy statement]
The portfolio exists to compound purchasing power over the long term while avoiding permanent capital impairment. Strategic allocation will remain within pre-defined bands. Broad, low-cost diversified holdings form the core. Direct securities may be held only when financial quality and conservative valuation establish a margin of safety. Market quotations are treated as transaction opportunities rather than as estimates of personal financial success. Speculation, if any, is separately capped and is not required to meet the plan's objectives.
\end{modernbox}

\section*{Primary Source}
Benjamin Graham, \textit{The Intelligent Investor: A Book of Practical Counsel}, revised edition, updated with commentary by Jason Zweig, preface and appendix material by Warren E. Buffett. Chapter references in this document point to the structure of that edition.

\end{document}
